%% ============================================================================
%% Fichier  : annexes/Ann_E_bibliographie.tex
%% Titre    : Bibliographie commentée
%% Version  : 1.0 -- novembre 2026
%% ============================================================================
%% Alignée sur les annexes des autres cours de l'auteur (ACSI, INF3226).
%% ============================================================================

\chapter{Bibliographie Commentée}
\label{ann:biblio}

Pour chaque partie, quelques lectures, avec ce qu'elles apportent. Les références
complètes figurent dans la bibliographie du livre. Les textes de l'auteur sont des
prépublications déposées sur l'IACR ePrint~; aucun travail non publié n'est cité.

\section*{Partie IV · Forensique opérationnelle}
\begin{itemize}
    \item Casey~\cite{casey2004} et Carrier~\cite{carrier2005}~: preuve numérique et systèmes de fichiers.
    \item NIST SP~800-86 et 800-101~\cite{nist80086,nist800101r1}~: intégrer la forensique à la réponse aux incidents~; les mobiles.
    \item ISO/IEC~27037~\cite{iso27037} et RFC~3227~\cite{rfc3227}~: collecte et préservation.
    \item Miller~\cite{miller1956} et Heuer~\cite{heuer1999}~: écrire un rapport lisible, éviter les biais.
\end{itemize}

\section*{Partie V · Défis émergents}
\begin{itemize}
    \item Les cadres de la Partie~IV, appliqués à l'anti-forensique, à l'IA et aux conteneurs~; les textes de l'Annexe~\ref{ann:lois} pour l'entraide.
\end{itemize}

\section*{Partie VI · Fondements cryptographiques}
\begin{itemize}
    \item Shannon~\cite{shannon1948} et Dolev et Yao~\cite{dolevyao1983}~: information et adversaire.
    \item Johnson et al.~\cite{johnson2001ecdsa}, Groth~\cite{groth2016zksnark}~: signatures et preuves à divulgation nulle.
    \item Adida~\cite{adida2009helios,adida2009university}~: un protocole vérifiable mis à l'épreuve.
    \item Le Trilemme CRO~\cite{minka2025cro} et ZK-NR~\cite{minka2025zknr}, prépublications de l'auteur sur l'IACR ePrint.
\end{itemize}
